\section{Experimental Setup}
\label{sec:setup}

\subsection{Data and Metric}
ERA5~\cite{hersbach2020} is represented on a regular $0.25^\circ$
latitude--longitude grid and cropped to $25$--$49^\circ$N and
$130$--$70^\circ$W ($97\times241$ cells). We use the
20 FourCastNet variables~\cite{pathak2022}: five surface fields; winds and
geopotential at 1000/850\,hPa; temperature and relative humidity at 850\,hPa;
wind, geopotential, temperature, and relative humidity at 500\,hPa; geopotential
at 50\,hPa; and total column water vapour. Frames are six-hourly at 00/06/12/18
UTC. Training uses 2014--2016 (4{,}384 frames), validation 2017 (1{,}460), and
evaluation 2018 (1{,}460). Architecture and codec parameters are selected using
validation data; reported quantisation steps are fixed before 2018 evaluation.
The full-year 2022 test uses all 1{,}460 timestamps, January 1 at 00 UTC through
December 31 at 18 UTC. ERA5 truth combines 18 native
ARCO-ERA5 fields with native NCAR relative humidity at 500/850\,hPa; decoded
CRA5 and truth timestamps, coordinates, and channel order are checked exactly.
Source packing precision is cross-checked against 2018 samples. The original
checkpoint and 2014--2016 statistics are frozen; neither test year is used for
fine-tuning or selection.

We compute channel nRMSE as $\sqrt{\langle(x_j-\hat x_j)^2\rangle}/\sigma_j$.
The mean pools test times and cells; $\sigma_j$ is the training standard deviation. We
average these channel errors equally. Primary CRA5 means use 19 direct channels;
the full 20-channel mean is labelled separately. CRA5 does not store total column
water vapour (\texttt{tcwv}); we derive it by pressure integration of specific
humidity, truncated at the surface. Its $-0.63\,\mathrm{kg\,m^{-2}}$ bias can
be learned by the corrector, so its gains partly repair this derivation.

For channel errors $b_j$ before and $a_j$ after correction, the headline reduction
is $G=1-\sum_j a_j/\sum_j b_j$. Equal-weight fractional gain is
$\bar G=p^{-1}\sum_j(1-a_j/b_j)$. The former weights fractional gains by baseline
error; the latter treats channels equally. Terrain roughness uses $G$; the legacy proxy-elevation control instead uses
$G_{\mathrm{pool}}$ (Section~\ref{sec:terrain}).

\subsection{Correction Network}
\label{sec:cnn}
All correctors train with 20 input/output channels; direct19 is only an evaluation
subset. U-Net depth is 2, base/bottleneck widths 64/128, with $3\times3$
convolutions and 1{,}646{,}548 parameters. Input is $(\hat x-\mu)/\sigma$,
target $(x-\hat x)/\sigma$, and corrected field $\hat x+\sigma f(\cdot)$;
$\mu,\sigma$ are fixed training statistics. CRA5 correction minimises normalised
residual MSE using AdamW: learning rate $3\times10^{-4}$, weight decay $10^{-4}$,
gradient clipping 1.0, batch 8, mixed precision, seed 0, maximum 200 epochs,
and early-stopping patience 12. The architecture comes from a 33-configuration
MPCA sweep ranked on 2017, not a CRA5 architecture search. Classical checkpoints
share the architecture/objective but not identical budgets: inspected MPCA-only
and e995 logs allow 200 and 100 epochs. Each is reused without retraining here.

We decode the released CRA5 streams for 2014--2018 and extract CONUS.
Baseline 20-channel nRMSE ranges from 0.08523 to 0.08752 across these years;
they overlap CRA5 training (Section~\ref{sec:limitations}). New 2018 diagnostics
replay the original checkpoint with float32 inference on the existing float16
normalised cache and float64 aggregation, reproducing original rounded results.

\subsubsection{Temporal uncertainty.}
Squared-error sums and counts for each frame and channel support a paired circular
moving-block bootstrap: 28-frame (seven-day) blocks wrap at the year boundary,
with each draw truncated to 1{,}460 frames. In 1{,}000 draws (seed 20261026),
before/after errors share sampled timestamps; channel RMSEs and $G$ are recomputed.
Percentile 95\% intervals condition on checkpoint, year, and grid, excluding
training-seed and spatial uncertainty. The same protocol is applied to 2022.
